%!TEX program = uplatex
\RequirePackage{plautopatch}
\documentclass[uplatex,dvipdfmx]{jlreq}
\usepackage{amsmath,amssymb}

\pagestyle{empty}

\begin{document}

%======================================================================
% 表紙
%======================================================================
\begin{center}
  \vspace*{20pt}
  {\Huge\bfseries ENCOUNTER with MATHEMATICS}\\[10pt]
  {\Large 第40回}\\[6pt]
  {\large 2007年5月18日、19日}\\[20pt]
  {\Huge\bfseries 力学系のゼータ関数}\\[6pt]
  {\Large -- 古典力学と量子力学のカオス --}\\[20pt]
  梗\qquad 概
\end{center}

今回の講演会では力学系のゼータ関数についていくつかの話題を取りあげたい。
力学系のゼータ関数は与えられた流れについて複素変数$s$を含む次の（無限）積として定義される：
\[
Z(s)=\prod_{k=0}^{\infty}\zeta(s+k)^{-1},\qquad\mbox{ただし} \quad  \zeta(s)=\prod_{\gamma}(1-e^{-s|\gamma|})^{-1}.
\]
ここで積$\prod_\gamma$は閉軌道$\gamma$についての無限積であり、$|\gamma|$は閉軌道$\gamma$の周期を表す。
このような関数は Smale によってSelberg ゼータ関数の力学系の言葉を用いた一般化として導入され、一方で既に古典となったアノソフ力学系について新たな精密な理論の構築を促し、一方では量子カオスなどの話題と関係するなど力学系の研究に奥行きと広がりを与えている。

\vfill

\newpage

%======================================================================
% 講演1：辻井正人「力学系のゼータ関数と Ruelle 転移作用素のスペクトル」
%======================================================================
\begin{center}
{\large\bfseries 力学系のゼータ関数と Ruelle 転移作用素のスペクトル}\\[6pt]
\hspace*{\fill}{\large\bfseries 辻井　正人（九州大学・数理学研究院）}
\end{center}

負曲率曲面上の測地流はアノソフ流の典型的な例である。特に負定曲率曲面の場合は Selberg の古典的な結果から測地流に対するゼータ関数は複素平面全体に正則関数として解析接続され、（非自明な）零点は曲面上のラプラス作用素の固有値を使って表されるという著しい事実が成り立つ。この事実は測地流の閉軌道という古典力学的な対象とラプラス作用素の固有値という量子力学的な対象を関係させているという点で示唆的である。ただ、これが「真に」示唆的であるとすると（そうである保証は何もないのだが）より一般の力学系（たとえば曲率が一定でない場合や一般のアノソフ流）について適当な形に一般化されると考えるのが自然であろう。これが力学系のゼータ関数を研究する一つの動機である。

本講演では、以上の研究の背景について簡単に説明した後、力学系のゼータ関数の解析的な性質を転移作用素のスペクトル的な性質と関係させる研究について解説する。このような研究は主に Ruelle によって進められてきたものであり、アイデアとしては、力学系のゼータ関数の解析的な性質を転移作用素のトレースの評価に帰着し、ゼータ関数の零点や極を転位作用素またはその生成作用素の固有値と関係させるというものである。離散力学系（拡大写像やアノソフ微分同相）の場合にはゼータ関数は上で述べたのとは多少違った形で定義されるが、その場合にはこのアイデアはかなりの部分実現されている。特に転移作用素のスペクトルについての議論は離散的な固有値は力学系の中に何らかの「共鳴」（Ruelle 共鳴）が起こっていることを示していて、それ自体興味深い。一方、本来の目標である流れの場合は現在も技術的に難しい部分が多いが、期待される結果を含めて話をしたい。

\vspace{5mm}\noindent
参考文献：V. Baladi, \textit{Positive transfer operators and decay of correlations.} World Scientific (2000)

\newpage

%======================================================================
% 講演2：首藤啓「量子カオスの諸問題と半古典ゼータ関数」
%======================================================================
\begin{center}
{\large\bfseries 量子カオスの諸問題と半古典ゼータ関数}\\[6pt]
\hspace*{\fill}{\large\bfseries 首藤　啓（首都大学東京・理工学研究科・物理学専攻）}
\end{center}

量子力学の時間発展を記述するシュレディンガー方程式は線形であることから、量子力学には古典力学と同じ意味でのカオスは存在しない。しかしながら、古典力学がカオスを示すと、対応する量子系には古典カオスの明確な「痕跡」が現れる。たとえば、完全可積分な系では、保存量によって定まるトーラス上に量子力学の固有関数が局在するのに対して、カオス系では、固有関数の台はエネルギー一定面全体に広がる。固有エネルギーに関しても、完全可積分な系でのエネルギーの分布はポアッソン分布に従うのに対して、カオス系ではウィグナー分布が現れること（Bohigas--Giannoni--Schmit の予想）などが知られている。

古典カオスの痕跡がなぜ量子力学に現れるのか？　この問いに答えるためには、古典と量子の間の橋渡しの言語が必要となる。ポアッソンの和公式、セルバーグの跡公式などは、それぞれ調和振動子、定負曲率面上の測地流に対して、古典系の周期軌道と量子系のスペクトルとの双対的な関係を「厳密に」表現するものとして、ここでの目的に適うもっとも良いお手本といえる。一方、より一般の系に対しては、プランク定数ゼロの極限（半古典極限）で漸近的に成り立つ類似の公式（カオス系の場合、グッツウィラーの跡公式）が知られており、数学的なレベルでの解析はほとんど進んでいないものの、固有エネルギーの分布に見られる普遍則の説明などには必須の道具となっている。さらに、定負曲率面上の測地流の問題におけるセルバーグのゼータ関数に対応して、グッツウィラー--ボロスのゼータ関数（半古典ゼータ関数）が半古典極限において導入される。力学系のゼータ関数同様、半古典ゼータ関数の解析的性質を明かにすることは量子カオスの重要課題のひとつであり、他のゼータ関数で知られている事実を手がかりにいくつかの試みがある。

講演では、跡公式、半古典ゼータ関数を用いた量子カオス系へのアプローチを紹介するとともに、量子カオスをめぐるいくつかの最近の話題についても触れる予定である。

\newpage

%======================================================================
% 講演3：盛田健彦「マルコフ系のゼータ関数」
%======================================================================
\begin{center}
{\large\bfseries マルコフ系のゼータ関数}\\[6pt]
\hspace*{\fill}{\large\bfseries 盛田　健彦（広島大学・大学院理学研究科）}
\end{center}

力学系ゼータ関数の解析的性質を調べる手段の一つとして転送作用素を用いた熱力学形式の方法がある。ここでは縮小写像の族が定める Markov 系の転送作用素と周期軌道分布に対応したゼータ関数との関わりについての入門的な話をしようと思う。おおまかな枠組みと一般論を紹介した後に、その応用例として次の2つの話題に焦点をしぼり話を進める予定である。

\medskip
(1)\ $\mathbb{H}$ は Poincaré 計量を備えた複素上半平面とし、
$PSL(2,\mathbb{Z})$ は、行列式が $1$ となる $2\times2$ 整数行列が定める1次分数変換の群とする。モジュラー曲面 $M_{1}=\mathbb{H}/PSL(2,\mathbb{Z})$ の Selberg ゼータ関数 $Z(s)$ は、連分数変換 $T_{G}\,:\, (0,1)\to (0,1)$ の2回反復合成 $T_{G}^{2}$ に附随する転送作用素の族 $L(s)$ の Fredholm 行列式で表現することができる。すなわち、$Z(s)=\mathrm{Det}(I-L(s))$ が成り立つ。これを利用すると Selberg の跡公式の助けをかりずにモジュラー曲面の閉測地線に関する素数定理型定理を導くことが可能となる。

\medskip
(2)\ 平面 $\mathbb{R}^{2}$ 内に正の曲率をもった単純閉曲線を境界とする障害物
$Q_{1},\,\dots,\, Q_{n}$ $(n\ge3)$ が、どの2つの凸包を考えてもその2つ以外の障害物と共有点をもたないように配置されていたとする。これらの障害物の外部領域 $Q=\mathbb{R}^{2}\setminus\bigcup_{j=1}^{n}\bar{Q_{j}}$ の内部では単位速度で並進運動し、境界 $\partial Q=\bigcup_{j=1}^{n}\partial Q_{j}$ では幾何光学の反射法則に従う質点の運動を記述する散乱開撞球系を考える。すなわち、境界付き多様体 $\bar{Q}$ 上の測地流で境界では完全弾性衝突するものを考える。このとき、この散乱開撞球系の周期軌道の分布に対応するゼータ関数 $\zeta_{Q}(s)$ は周期軌道に関する素数定理型定理を示すに十分な解析的性質をもっていることがわかる。

\medskip
さらに、時間が許せば、これらの例に関連した話題や一般化の方向についても若干触れてみたいと思う。

\end{document}
